% Rubinstein - Capablanca, San Sebastian 1911, with Capablanca's notes.
% Build with: python3 chessbook.py rubinstein-capablanca.tex
\documentclass[10pt,twoside]{article}
\usepackage[paperwidth=487.56bp,paperheight=691.65bp,top=72pt,bottom=44pt,inner=50pt,outer=48pt,
            headheight=13pt,headsep=20pt,footskip=0pt]{geometry}
\usepackage[T1]{fontenc}
\usepackage[lining]{ebgaramond}
\usepackage{graphicx}
\usepackage{microtype}
\usepackage{tikz}
\usepackage{fancyhdr}
\usepackage[hidelinks,pdftitle={Rubinstein - Capablanca, San Sebastian 1911},
            pdfauthor={J. R. Capablanca}]{hyperref}

\newcommand{\textsizebook}{\fontsize{10.5}{12.6}\selectfont}
\AtBeginDocument{\textsizebook}
\setlength{\columnsep}{11pt}
\setlength{\parindent}{9pt}
\setlength{\parskip}{0pt}
\setlength{\emergencystretch}{1.5em}
\raggedbottom

\newlength{\boardwidth}
\setlength{\boardwidth}{160pt}

% A main line move, as a bold heading above its note
\newcommand{\Move}[1]{\par\addvspace{8pt}\noindent{\bfseries #1}\par\nopagebreak}
% A diagram in the text
\newcommand{\ChessDiagram}[1]{\par\addvspace{7pt}{\centering #1\par}\addvspace{5pt}}

\pagestyle{fancy}
\fancyhf{}
\renewcommand{\headrulewidth}{0pt}
\fancyhead[LE,RO]{\thepage}
\fancyhead[CE]{Queen's Gambit Declined}
\fancyhead[CO]{Rubinstein -- Capablanca}

\game{d4 d5 Nf3 c5 c4 e6 cxd5 exd5 Nc3 Nc6 g3 Be6 Bg2 Be7 O-O Rc8 dxc5 Bxc5 Ng5 Nf6 Nxe6 fxe6
      Bh3 Qe7 Bg5 O-O Bxf6 Qxf6 Nxd5 Qh6 Kg2 Rcd8 Qc1 exd5 Qxc5 Qd2 Qb5 Nd4 Qd3 Qxd3 exd3 Rfe8
      Bg4 Rd6 Rfe1 Rxe1 Rxe1 Rb6 Re5 Rxb2 Rxd5 Nc6 Be6+ Kf8 Rf5+ Ke8 Bf7+ Kd7 Bc4 a6 Rf7+ Kd6
      Rxg7 b5 Bg8 a5 Rxh7 a4 h4 b4 Rh6+ Kc5 Rh5+ Kb6 Bd5 b3 axb3 a3 Bxc6 Rxb3 Bd5 a2 Rh6+}

\begin{document}

% ---- Opening page ----
\thispagestyle{empty}
\begin{tikzpicture}[remember picture,overlay]
  \node[circle,draw=black,line width=1.6pt,fill=black!20,minimum size=112pt] (seal)
       at ([xshift=75pt,yshift=-75pt]current page.north west) {};
  \draw[line width=0.9pt] ([shift={(-48pt,-30pt)}]seal.center) -- ([shift={(44pt,34pt)}]seal.center);
  \node[rotate=37] at ([shift={(-11pt,14pt)}]seal.center) {\fontsize{25}{25}\selectfont\bfseries Game};
  \node[rotate=37] at ([shift={(10pt,-16pt)}]seal.center) {\fontsize{38}{38}\selectfont\bfseries 2};
  \node[anchor=north east] at ([xshift=-14pt,yshift=-24pt]current page.north east)
       {\board[after=2...c5,width=108pt,bare]};
\end{tikzpicture}

\vspace*{38pt}
\begin{center}
  \hspace*{56pt}\begin{minipage}{235pt}\centering
    {\fontsize{22}{25}\selectfont\bfseries Queen's Gambit\\Declined\par}
    \vspace{4pt}\rule{\linewidth}{1pt}\\[-9.5pt]\rule{\linewidth}{0.4pt}
  \end{minipage}\hspace*{56pt}

  \vspace{14pt}
  {\fontsize{19}{22}\selectfont\bfseries Rubinstein -- Capablanca\par}
  \vspace{12pt}
  {\fontsize{15}{18}\selectfont\bfseries San Sebastian 1911\par}
  \vspace{16pt}
  {\fontsize{11.5}{14}\selectfont\bfseries
   \moves{1.d4 d5 2.Nf3 c5 3.c4 e6 4.cxd5 exd5}\par\vspace{2pt}
   \rule{300pt}{0.6pt}\par\vspace{2pt}
   \moves{5.Nc3 Nc6 6.g3 Be6 7.Bg2 Be7 8.O-O Rc8}\par}
\end{center}

\vspace{-2pt}
{\setlength{\parindent}{0pt}\fontsize{11.5}{14.4}\selectfont\bfseries
\moves{6...Be6} avoiding the beaten track\hfill\hyperref[be6]{\pageref*{be6}}\par
\moves{8...Rc8} a change with poor success\hfill\hyperref[rc8]{\pageref*{rc8}}\par
\moves{14.Bxf6} Rubinstein's combination begins\hfill\hyperref[bxf6]{\pageref*{bxf6}}\par
\moves{16.Kg2!} the move Black had not considered\hfill\hyperref[kg2]{\pageref*{kg2}}\par
\moves{30...a6} a lost chance to draw\hfill\hyperref[asix]{\pageref*{asix}}\par
\moves{38.Bd5} a blunder that goes unpunished\hfill\hyperref[bd5]{\pageref*{bd5}}\par}

\vfill
\noindent\begin{minipage}[t]{\linewidth}
\centering\small
\begin{tabular}{@{}c@{\hspace{13pt}}c@{\hspace{13pt}}c@{}}
  after \moves{14...Qxf6} & after \moves{15...Qh6} & after \moves{37...Kb6} \\[2pt]
  \board[after=14...Qxf6,width=114pt,bare] & \board[after=15...Qh6,width=114pt,bare] &
  \board[after=37...Kb6,width=114pt,bare] \\[3pt]
  \moves{15.Nxd5!} & \moves{16.Kg2!} & \moves{38.Bd5}
\end{tabular}
\end{minipage}

\vspace{13pt}
{\centering\fontsize{8.5}{10}\selectfont\itshape
 Notes by J. R. Capablanca, from \textup{Chess Fundamentals} (1921), which is in the public domain.\\
 Descriptive notation has been converted to algebraic. Click any move or diagram to open it on Lichess.\par}
\clearpage

% ---- The game ----
\twocolumn
\textsizebook

\mainline{1.d4 d5 2.Nf3 c5 3.c4 e6 4.cxd5 exd5 5.Nc3 Nc6 6.g3 Be6}\label{be6}

\moves{6...Nf6} is the normal move in this variation. White's development was first introduced by
Schlechter and elaborated later on by Rubinstein. It aims at the isolation of Black's d-pawn,
against which the White pieces are gradually concentrated. In making the text move I was trying to
avoid the beaten track. Being a developing move there should be no objection to it in the way of
general principles, except that the Knights ought to come out before the Bishops.

\mainline{7.Bg2 Be7 8.O-O Rc8}\label{rc8}

In pursuance of the idea of changing the normal course of this variation, but with very poor
success. The move in theory ought to be unsound, since Black's king's knight is yet undeveloped. I
had not yet learned of the attack founded on \figN g5 and the exchange of the Bishop on e6. Either
\moves{8...Nf6} or \moves{8...h6}, to prevent either \figB{} or \figN g5, was right.

\board

\mainline{9.dxc5 Bxc5 10.Ng5 Nf6 11.Nxe6 fxe6 12.Bh3 Qe7 13.Bg5 O-O}

This is a mistake. The right move was \moves{13...Rd8} in order to get the Rook away from the line
of the Bishop on h3 and at the same time to support the d-pawn. Incidentally it shows that White
failed to take proper advantage of Black's weak opening moves. Against the text move White makes a
very fine combination which I had seen, but which I thought could be defeated.

\mainline{14.Bxf6 Qxf6}\label{bxf6}

I considered \moves{14...gxf6}, which it seemed would give me a playable game, but I thought White's
combination unsound and therefore let him play it, to my lasting regret.

\board

\mainline{15.Nxd5! Qh6}

\board

\mainline{16.Kg2!}\label{kg2}

\emph{This} is the move which I had not considered. I thought that Rubinstein would have to play
\moves{16.Bg2}, when I had in mind the following winning combination:
\moves[as=a]{16.Bg2 Ne5! 17.Nf4} (if \moves[from=a]{17.Rc1 Qxc1!! 18.Qxc1 Bxf2+} wins),
\moves[from=a,as=b]{17...Ng4 18.h3} (if \moves[from=b]{18.Nh3 Bxf2+} wins the exchange),
\moves[from=b]{18...Nxf2 19.Rxf2 Bxf2+ 20.Kxf2 g5}, and Black should win. It is curious that this
combination has been overlooked. It has been taken for granted that I did not see the 17th move
\moves{17.Qc1}.

\mainline{16...Rcd8}

After White's last move there was nothing for me to do but submit to the inevitable.

\mainline{17.Qc1! exd5 18.Qxc5 Qd2 19.Qb5 Nd4 20.Qd3 Qxd3 21.exd3 Rfe8 22.Bg4}

This gives Black a chance. He should have played \moves[as=c]{22.Rfe1}. If then
\moves[from=c,as=d]{22...Nc2 23.Rxe8+ Rxe8 24.Rc1 Re2 25.Kf1}
\moves[from=d,as=e]{25...Nd4} (if \moves[from=d]{25...Rd2 26.Be6+ Kf8 27.Bxd5} would win);
\moves[from=e]{26.Rc8+ Kf7 27.Rc7+ Re7 28.Rc5} wins.

\mainline{22...Rd6 23.Rfe1 Rxe1 24.Rxe1 Rb6 25.Re5 Rxb2 26.Rxd5 Nc6 27.Be6+ Kf8 28.Rf5+ Ke8
          29.Bf7+ Kd7 30.Bc4}

\board

\mainline{30...a6}\label{asix}

A bad move, which gives away any legitimate chance Black had to draw. It loses a very important
move. In fact, as the course of the game will show, it loses several moves. The proper way was to
play \moves[as=f]{30...Kd6}. If then \moves[from=f]{31.Rb5 Rxb5 32.Bxb5 Nd4}, followed by ...b5;
and White would have an exceedingly difficult game to draw on account of the dominating position of
the Knight on d4 in conjunction with the extra Pawn on the Queen's side and the awkward position of
White's King. (See how this is so.)

\mainline{31.Rf7+ Kd6 32.Rxg7 b5 33.Bg8 a5 34.Rxh7 a4 35.h4 b4 36.Rh6+ Kc5 37.Rh5+ Kb6 38.Bd5}\label{bd5}

With these last three moves White again gives Black a chance. Even before the last move
\moves{38.Bc4} would have won with comparative ease, but the text move is a downright blunder, of
which, fortunately for him, Black does not avail himself.

\board

\mainline{38...b3}

\moves[as=g]{38...Rxa2} would make it practically impossible for White to win, if he can win at all.
White's best continuation then would have been:
\moves[from=g,as=h]{39.Bc4 Rc2 40.Rb5+ Kc7 41.Bg8 a3 42.h5 a2 43.Bxa2 Rxa2}, and if there is a win it
is very difficult to find it, as against \moves[from=h]{44.h6 Ra3!} offers excellent chances for a
draw.

\mainline{39.axb3 a3 40.Bxc6 Rxb3}

If \moves{40...a2 41.Rb5+ Ka6 42.Rb8}.

\mainline{41.Bd5 a2 42.Rh6+}

\noindent\textbf{Resigns.}

As an end game, this is rather a sad exhibition for two masters. The redeeming feature of the game
is Rubinstein's fine combination in the middle game, beginning with \moves{14.Bxf6}.

\end{document}
